% ====================================================================
% FILE: content/08_boundary_ru.tex
% Расширенная граница и геометрия патчей — Core 1.3.3
% ====================================================================

\section{Расширенная граница и геометрия патчей}
\label{sec:boundary-extended}

Эта глава восстанавливает полную теорию границ и геометрии патчей.
Границы составляют одну из самых универсальных структур в рамках OC:
они управляют фазовым разделением в \(K_2\), формированием мембран в \(K_3\)--\(K_4\),
возбудимостью в \(K_5\), а также представительными или институциональными пределами в
\(K_6\)--\(K_8\).  
Текущая версия консолидирует полный набор конструкций, необходимый для Core~1.3.3:
пороги, задаваемые границей, дискретизация патчей, локальные классы порогов
и переходы, инициируемые границей.

\subsection{Формальное определение \texorpdfstring{$\partial\Omega(K)$}{границы}}

Для любого континуума \(K\) с допустимым пространством состояний \(\Omega(K)\) и
пороговым ландшафтом \(\Theta(K)\) граница определяется как множество состояний, на которых
по крайней мере один порог достигает точного насыщения:
\[
  \partial\Omega(K)
  =
  \big\{\, s \in \overline{\Omega(K)} :
          \exists k\ \text{с таким, что}\ g_k(s) = \Theta_k \big\}.
\]
Эквивалентно, если \(f_k(s)=g_k(s)-\Theta_k\) обозначает сдвинутую остаточную форму,
используемую в приложении обозначений и главе модели, то внутренние состояния
удовлетворяют \(f_k(s)<0\), состояния допустимой стороны — \(f_k(s)\le0\),
а эта же граница есть нулевое множество \(f_k(s)=0\) для хотя бы одного
активного порога.

Итак:
\begin{itemize}
    \item внутренние состояния удовлетворяют \(f_k(s) < \Theta_k\) для всех порогов;
    \item внешние состояния нарушают по крайней мере один порог;
    \item граница — это \emph{поверхность, определяемая порогом}, а не геометрический
          примитив.
\end{itemize}

Эта формулировка объединяет классические поверхности, мембраны, везикулы,
перколяционные фронты, возбудимые границы, логические пределы и институциональные
пороги в единой структуре.

\subsection{Геометрия границы}

Геометрия границы возникает из локальной структуры порогов, потенциалов и потоков.
Если активный порог в точке \(s \in \partial\Omega(K)\) есть \(g_k\), то внешняя
нормаль ориентирована по градиенту:
\[
  n(s) \propto \nabla g_k(s).
\]

Несколько активных порогов создают углы, ребра и области высокой кривизны.
Кривизна становится динамически значимой всякий раз, когда в потенциал входят
члены изгиба или поверхностной энергии:
\[
  P_{\mathrm{bend}}(s) = \tfrac{1}{2}\,\kappa(s)^2\,\kappa_b,
\]
с жесткостью изгиба \(\kappa_b\).

Это связывает:
\begin{itemize}
    \item \(K_2\): классическое поверхностное натяжение,
    \item \(K_3\)--\(K_4\): кривизна мембраны и энергетика везикул,
    \item \(K_5\): геометрию возбудимых границ.
\end{itemize}

\subsection{Модель патчей для границ}

OC рассматривает границы не как непрерывные однородные поверхности, а как
наборы динамических \emph{патчей}:
\[
  \partial\Omega(K) = \bigcup_i \partial\Omega_i.
\]

Каждый патч \(P_i\) несет:
\[
  (\sigma_i,\Theta_i,P_i(t),J_i(t)),
\]
в котором отражаются его локальное состояние, пороги, потенциалы и потоки.

\paragraph{Локальные состояния.}
Патчи принимают символические состояния:
\[
  \sigma_i \in \{L_\alpha, L_\beta, L_o, L_f, L_b\},
\]
кодирующие локальные упаковочные или структурные режимы:
\begin{itemize}
    \item \(L_\alpha\): жидко-неупорядоченный,
    \item \(L_\beta\): упорядоченный/гелевый,
    \item \(L_o\): жидко-упорядоченный,
    \item \(L_f\): хрупкий/тонкий (предразрывной),
    \item \(L_b\): разрушенное граничное состояние.
\end{itemize}

Этот алфавит обобщается за пределы биологических мембран до перколяционных краев,
линий разрушения, флуктуирующих интерфейсов и даже символических/когнитивных
границ.

\paragraph{Локальные пороги.}
Каждый патч имеет вектор:
\[
  \Theta_i = (\Theta^{\rm perm}_i, \Theta^{\rm grad}_i, \Theta^{\rm mem}_i,
              \Theta^{\rm bend}_i, \ldots ),
\]
с соответствующими локальными потенциалами \(P_i(t)\) и потоками \(J_i(t)\).

\subsection{Система порогов границы}

Для диапазона \(K_2\)--\(K_5\) универсально возникают три класса порогов.

\paragraph{1. Порог проницаемости \(\Theta_{\rm perm}\).}
Патч становится проницаемым, когда:
\[
  |\nabla P_i| > \Theta^{\rm perm}_i.
\]
Примеры:
\begin{itemize}
    \item осмотическая проницаемость (\(K_3\)--\(K_4\)),
    \item начало ионного утечки (\(K_5\)),
    \item проникновение фазового фронта (\(K_2\)).
\end{itemize}

\paragraph{2. Порог градиента \(\Theta_{\rm grad}\).}
Если градиент неустойчив:
\[
  |\nabla P_i| > \Theta^{\rm grad}_i
  \quad\Rightarrow\quad \sigma_i \to L_f.
\]
Это объясняет:
\begin{itemize}
    \item разрыв везикул при высоком осмотическом давлении (\(K_4\)),
    \item инициацию потенциала действия в начальных сегментах аксона (\(K_5\)),
    \item ускоренные фронты фаз и распространение трещин (\(K_2\)).
\end{itemize}

\paragraph{3. Порог целостности мембраны \(\Theta_{\rm mem}\).}
Если натяжение мембраны или кривизна превышают пределы:
\[
  \gamma_i > \Theta^{\rm mem}_i
  \qquad\text{или}\qquad
  \kappa(s) > \Theta^{\rm bend}_i,
\]
то:
\[
  \sigma_i \to L_b.
\]

Это дает единое понимание разрыва, потери устойчивости, образования пор,
разрушения границы и коллапса поверхности.

\subsection{Динамика, задаваемая границей}

Соседние патчи взаимодействуют через потоки:
\[
  J_{ij} = F(P_i, P_j, \sigma_i, \sigma_j).
\]

Это приводит к:
\begin{itemize}
    \item стабилизации и перестройке мозаики патчей,
    \item распространению фронтов разрыва,
    \item миграции, управляемой кривизной,
    \item запуску пор,
    \item поведению возбудимой среды в \(K_5\).
\end{itemize}

Эволюция патчей описывается:
\[
  \frac{d\sigma_i}{dt}
  = S(\sigma_i, P_i, J_i, \Theta_i),
\]
где \(S\) — оператор структурного обновления, не зависящий от области.

\medskip
\textbf{Принцип локализации границы.}
\emph{Все переходы размерностей \(K_3\)--\(K_5\) запускаются на граничных патчах.}

Это объясняет:
\begin{itemize}
    \item разрушение протоклеток, начинающееся в точках высокой кривизны,
    \item запуск потенциалов действия в начальных сегментах аксона,
    \item нуклеацию фазовых переходов на дефектах или поверхностях.
\end{itemize}

\subsection{Провал и коллапс границы}

Коллапс границы происходит, когда связанный набор патчей становится разрушенным:
\[
  \sigma_i = L_b
  \qquad \forall\, i \in \mathcal{C},
\]
и кластер \(\mathcal{C}\) перколирует через всю границу.

По Теореме~3 (раздел Результаты):
\[
  \partial\Omega(K)\ \text{коллапсирует}
  \quad\Longrightarrow\quad
  \Omega(K) = \emptyset,
\]
и континуум гибнет.

Мы различаем:
\begin{itemize}
    \item \textbf{локальный отказ}: изолированные патчи \(L_b\), обратимые;
    \item \textbf{коллапс кластера}: перколирующая сеть \(L_b\), необратимая;
    \item \textbf{глобальная потеря границы}: исчезновение границы, влекущее
          смерть или переход в новый континуум.
\end{itemize}

\subsection{Граница и возрождение}

Рождение нового континуума \(K_{x+1}\) обычно требует:
\begin{itemize}
    \item формирования или разрыва граничных патчей,
    \item создания новых типов патчей,
    \item появления новых классов порогов в граничных областях,
    \item установления новой глобальной геометрии границы.
\end{itemize}

\textbf{Принцип возрождения границы.}
\emph{Новые континуумы возникают из новых геометрий границ.}

Примеры:
\begin{itemize}
    \item замыкание мембраны при переходе \(K_3 \to K_4\),
    \item появление возбудимых границ при переходе \(K_4 \to K_5\),
    \item формирование институциональных границ при переходе \(K_6 \to K_7\),
    \item возникновение представительных границ при переходе \(K_9 \to K_{10}\).
\end{itemize}

Следовательно, границы — не просто пределы:
они — генеративные поверхности, на которых возникают новые оси, ограничения и
циклы.
